\chapter{Introducci\'on}
\label{ch:introduction}

La inteligencia artificial ha experimentado un crecimiento exponencial en los
\'ultimos a\~nos, no solo como campo de investigaci\'on cient\'ifica sino como
una herramienta que se introduce r\'apidamente en la vida cotidiana de las
personas. La tendencia es inequ\'ivoca, y es posible que nos encontremos ante
uno de los cambios m\'as significativos en la historia de la tecnolog\'ia.

La historia del aprendizaje autom\'atico se remonta a los a\~nos 1940--1950, con
los primeros intentos de construir modelos computacionales inspirados en el
cerebro. El perceptr\'on, propuesto por Frank Rosenblatt en
1958~\citep{rosenblatt1958perceptron}, fue uno de los primeros algoritmos de
aprendizaje capaces de clasificar patrones linealmente separables. Sus
limitaciones te\'oricas, precisadas por Minsky y Papert~\citep{minsky1969perceptrons},
contribuyeron al primer ``invierno de la inteligencia artificial'': un periodo
prolongado de escepticismo y reducci\'on de la financiaci\'on. El resurgimiento
lleg\'o en los a\~nos 1980 con el algoritmo de retropropagaci\'on
(\emph{backpropagation})~\citep{rumelhart1986backprop}, que permiti\'o entrenar
redes multicapa de forma eficiente. Aun as\'i, el desvanecimiento del gradiente
y la limitada capacidad de c\'omputo de la \'epoca mantuvieron estas t\'ecnicas
en un segundo plano durante a\~nos.

La transformaci\'on que hoy conocemos como \emph{aprendizaje profundo}
(\emph{deep learning}) lleg\'o en 2012, cuando AlexNet~\citep{krizhevsky2012alexnet}
gan\'o la competici\'on ImageNet y redujo dr\'asticamente el error de
clasificaci\'on de los m\'etodos tradicionales. Este \'exito se apoy\'o en la
convergencia de tres factores: grandes conjuntos de datos etiquetados, hardware
especializado accesible (GPU) capaz de entrenar modelos con millones de
par\'ametros, y nuevas t\'ecnicas de entrenamiento que superaban los obst\'aculos
hist\'oricos de las redes profundas. Desde entonces, el aprendizaje profundo ha
demostrado una eficacia notable en reconocimiento de im\'agenes, procesamiento de
lenguaje natural, generaci\'on de contenido y predicci\'on de estructuras
prote\'icas, entre muchas otras tareas.

Las arquitecturas responsables de estos logros ---redes convolucionales,
recurrentes y de atenci\'on--- fueron concebidas para operar sobre datos
organizados en dominios euclidianos regulares: rejillas de p\'ixeles, secuencias
temporales, tablas de caracter\'isticas. A medida que la comunidad comenz\'o a
aplicarlas a datos con estructura intr\'insecamente m\'as rica (mol\'eculas, redes
sociales, superficies tridimensionales), se hizo patente que la suposici\'on de
regularidad euclidiana constitu\'ia una limitaci\'on fundamental ---conceptual, y
no meramente t\'ecnica---. Esta constataci\'on impuls\'o un nuevo paradigma, el
\emph{aprendizaje profundo geom\'etrico} (\emph{geometric deep learning}, GDL),
que invierte la relaci\'on tradicional entre datos y arquitectura: en lugar de
forzar los datos a encajar en modelos preexistentes, son las propiedades
geom\'etricas de los datos las que dictan la estructura del modelo.

\section{¿Qu\'e es el aprendizaje profundo geom\'etrico?}
\label{sec:what-is-gdl}

Las arquitecturas tradicionales est\'an dise\~nadas para datos con estructura
euclidiana regular. Las redes convolucionales aprovechan la rejilla regular de
una imagen, donde cada p\'ixel tiene un n\'umero fijo de vecinos en posiciones
predefinidas; los modelos recurrentes y de atenci\'on procesan secuencias dotadas
de un orden temporal natural. Esta dependencia de los dominios regulares se
rompe cuando los datos no se ajustan a tales plantillas.

Consid\'erese el an\'alisis de una red social: los usuarios y sus conexiones
forman un grafo irregular en el que cada nodo puede tener un n\'umero arbitrario
de vecinos, sin un orden natural ni una rejilla. Aplicar directamente una red
convolucional exigir\'ia aplanar el grafo en una representaci\'on euclidiana (como
una matriz de adyacencia), descartando informaci\'on estructural crucial e
introduciendo una gran cantidad de ceros que hacen el c\'omputo ineficiente. El
an\'alisis de mol\'eculas en el descubrimiento de f\'armacos es igual de
ilustrativo: las mol\'eculas son, de forma natural, grafos cuyos nodos son
\'atomos y cuyas aristas son enlaces qu\'imicos, y su comportamiento
farmacol\'ogico viene gobernado por la geometr\'ia tridimensional y la topolog\'ia
de los enlaces. Los m\'etodos cl\'asicos depend\'ian de descriptores dise\~nados a
mano o de codificaciones de cadena (como SMILES), perdiendo informaci\'on
geom\'etrica esencial para la actividad biol\'ogica. La misma ``maldici\'on de la no
euclidianidad'' aparece en el an\'alisis de formas 3D, donde las superficies no
admiten una parametrizaci\'on natural sobre una rejilla regular y t\'ecnicas como
la voxelizaci\'on o la proyecci\'on a m\'ultiples vistas introducen artefactos y
borran detalle geom\'etrico fino.

La \Cref{fig:data-types} ilustra las tres categor\'ias fundamentales de datos
seg\'un su estructura geom\'etrica: datos euclidianos con rejilla regular, grafos
con topolog\'ia irregular y variedades con geometr\'ia continua.

\begin{figure}[h]
\centering
\begin{tikzpicture}[scale=0.9]
% --- Panel izquierdo: euclidiano (rejilla regular) ---
\begin{scope}[shift={(-5,0)}]
    \begin{scope}[on background layer]
        \fill[gdlblue!15, rounded corners=2pt] (-0.3,-0.3) rectangle (3.3,3.3);
    \end{scope}
    \foreach \x in {0,...,3} \foreach \y in {0,...,4}
        \draw[gdlgray!50] (\x*0.75, \y*0.75) -- (\x*0.75+0.75, \y*0.75);
    \foreach \x in {0,...,4} \foreach \y in {0,...,3}
        \draw[gdlgray!50] (\x*0.75, \y*0.75) -- (\x*0.75, \y*0.75+0.75);
    \foreach \x in {0,...,4} \foreach \y in {0,...,4}
        \fill[gdlblue!60] (\x*0.75, \y*0.75) circle (2.5pt);
    \node[below, font=\bfseries, text=gdlblue!70!black] at (1.5, -0.7) {Euclidiano};
\end{scope}
% --- Panel central: grafo (topolog\'ia irregular) ---
\begin{scope}[shift={(0.5,0)}]
    \coordinate (n1) at (0.3, 2.8); \coordinate (n2) at (1.8, 3.0);
    \coordinate (n3) at (0.0, 1.5); \coordinate (n4) at (1.5, 1.6);
    \coordinate (n5) at (2.8, 2.0); \coordinate (n6) at (0.5, 0.2);
    \coordinate (n7) at (2.2, 0.3);
    \foreach \a/\b in {n1/n2,n1/n3,n1/n4,n2/n4,n2/n5,n3/n4,n3/n6,n4/n5,n4/n6,n4/n7,n5/n7,n6/n7}
        \draw[gdlgreen!50, thick] (\a) -- (\b);
    \foreach \n in {n1,n2,n3,n4,n5,n6,n7} \fill[gdlgreen!70] (\n) circle (4pt);
    \node[below, font=\bfseries, text=gdlgreen!70!black] at (1.5, -0.7) {Grafo};
\end{scope}
% --- Panel derecho: variedad (geometr\'ia continua) ---
\begin{scope}[shift={(5.5,0)}]
    \shade[ball color=gdlorange!25, opacity=0.9]
        (0.0, 1.5) .. controls (-0.5, 2.8) and (1.0, 3.5) .. (1.8, 3.0)
        .. controls (2.6, 2.5) and (3.3, 2.8) .. (3.2, 1.8)
        .. controls (3.1, 0.8) and (2.5, -0.2) .. (1.5, 0.0)
        .. controls (0.5, 0.2) and (0.5, 0.3) .. (0.0, 1.5);
    \draw[thick, gdlorange!60]
        (0.0, 1.5) .. controls (-0.5, 2.8) and (1.0, 3.5) .. (1.8, 3.0)
        .. controls (2.6, 2.5) and (3.3, 2.8) .. (3.2, 1.8)
        .. controls (3.1, 0.8) and (2.5, -0.2) .. (1.5, 0.0)
        .. controls (0.5, 0.2) and (0.5, 0.3) .. (0.0, 1.5);
    \draw[gdlorange!30, thin] (0.5, 1.2) .. controls (1.2, 1.8) .. (2.5, 1.5);
    \draw[gdlorange!30, thin] (0.8, 2.2) .. controls (1.5, 2.6) .. (2.6, 2.2);
    \node[below, font=\bfseries, text=gdlorange!70!black] at (1.5, -0.7) {Variedad};
\end{scope}
\end{tikzpicture}
\caption[Tipos de datos geom\'etricos]{Los tres tipos fundamentales de datos
seg\'un su estructura geom\'etrica: datos euclidianos sobre una rejilla regular
(izquierda), grafos con topolog\'ia irregular donde cada nodo puede tener un
n\'umero distinto de vecinos (centro) y variedades con geometr\'ia continua que
requieren parametrizaciones locales (derecha).}
\label{fig:data-types}
\end{figure}

El aprendizaje profundo geom\'etrico es una respuesta matem\'aticamente rigurosa a
estas limitaciones. En lugar de adaptar los datos a las arquitecturas existentes
---proceso que invariablemente pierde informaci\'on---, dise\~na arquitecturas que
respetan y explotan la estructura geom\'etrica intr\'inseca de los datos. La
observaci\'on central es que muchos dominios del mundo real poseen simetr\'ias y
estructura geom\'etrica naturales que pueden aprovecharse para mejorar tanto la
eficiencia computacional como el rendimiento predictivo. Al incorporar estas
simetr\'ias directamente en la red, el GDL no solo resuelve los problemas de
representaci\'on anteriores, sino que tambi\'en garantiza propiedades deseables
como la equivarianza, la generalizaci\'on robusta y la eficiencia param\'etrica.

Los fundamentos matem\'aticos completos ---la teor\'ia de grupos, las nociones de
invarianza y equivarianza y el marco unificador de las ``5~Ges''--- se desarrollan
en los \Cref{ch:foundations,ch:priors}. Por ahora basta con entender que el GDL
proporciona un marco te\'orico que unifica arquitecturas aparentemente dispares
---redes convolucionales, redes de grafos y mecanismos de atenci\'on--- bajo un
paradigma geom\'etrico com\'un.

\section{La simetr\'ia como principio de dise\~no}
\label{sec:symmetry-principle}

La idea organizadora del aprendizaje profundo geom\'etrico es la
\emph{simetr\'ia}. Una simetr\'ia de un problema es una transformaci\'on de los
datos que deja invariable la magnitud de inter\'es. Rotar la imagen de un gato
sigue mostrando un gato; reetiquetar los nodos de un grafo molecular no cambia la
mol\'ecula. Tales transformaciones forman un \emph{grupo}, y exigir que un modelo
se comporte de forma coherente bajo ese grupo es una restricci\'on poderosa sobre
su arquitectura.

Dos nociones lo precisan. Una funci\'on es \emph{invariante} cuando su salida no
cambia al transformar la entrada mediante el grupo ---adecuada, por ejemplo, para
la clasificaci\'on de una imagen completa---. Es \emph{equivariante} cuando
transformar la entrada produce una salida transformada de forma correspondiente:
si $f$ es equivariante respecto de un grupo $\group$ que act\'ua sobre sus
espacios de entrada y salida, entonces
\[
  f(g \cdot x) = g \cdot f(x) \qquad \text{para todo } g \in \group .
\]
La equivarianza es la propiedad que deseamos en las capas intermedias de una red:
cada capa debe transformarse de forma predecible bajo la simetr\'ia, de modo que
la estructura se preserve a medida que la informaci\'on fluye por el modelo, con
una lectura invariante solo al final. Como la composici\'on de aplicaciones
equivariantes es de nuevo equivariante, este principio proporciona una receta
constructiva para arquitecturas completas, y no una propiedad que verificar a
posteriori. Estas ideas se formalizan en el \Cref{ch:priors}.

\section{Las 5 Ges: una hoja de ruta}
\label{sec:five-gs}

El aprendizaje profundo geom\'etrico organiza sus dominios y simetr\'ias en torno
a cinco temas recurrentes ---las ``5~Ges'' que dan subt\'itulo a este libro:
rejillas, grupos, grafos, geod\'esicas y gauges~\citep{bronstein2021geometric}---.
La \Cref{tab:five-gs} las resume y se\~nala el cap\'itulo en el que se desarrolla
cada una.

\begin{table}[h]
\centering
\caption{Las cinco Ges del aprendizaje profundo geom\'etrico.}
\label{tab:five-gs}
\begin{tabular}{@{}lp{6.2cm}p{3.6cm}@{}}
\toprule
\textbf{Tema} & \textbf{Idea} & \textbf{Cap\'itulo} \\
\midrule
Rejillas & Dominios euclidianos regulares con simetr\'ia de traslaci\'on & Redes convolucionales (\Cref{ch:grids}) \\
\addlinespace
Grupos & Simetr\'ias globales m\'as all\'a de la traslaci\'on, como rotaciones y reflexiones & Redes equivariantes a grupos (\Cref{ch:groups}) \\
\addlinespace
Grafos & Dominios relacionales irregulares con simetr\'ia de permutaci\'on & Redes neuronales de grafos (\Cref{ch:graphs}) \\
\addlinespace
Geod\'esicas & Geometr\'ia intr\'inseca de las variedades: distancias, curvatura y flujo de informaci\'on & Aprendizaje en variedades (\Cref{ch:geodesics}) \\
\addlinespace
Gauges & Marcos de referencia locales y la libertad en la elecci\'on de coordenadas & Redes equivariantes gauge (\Cref{ch:gauges}) \\
\bottomrule
\end{tabular}
\end{table}

Estos temas no est\'an aislados. Las rejillas son el dominio m\'as r\'igido y el
punto de partida hist\'orico; relajar la exigencia de una rejilla global
conservando las simetr\'ias globales conduce a las redes equivariantes a grupos;
descartar por completo la rejilla ambiente y conservar solo la estructura
relacional conduce a los grafos; dotar a un dominio de distancias y curvatura
intr\'insecas conduce a las variedades y las geod\'esicas; y reconocer que en un
dominio curvo no existe una elecci\'on can\'onica de marco local conduce a los
gauges. Un \'unico hilo ---el paso de mensajes--- recorre todos ellos y revela las
redes convolucionales, las redes de grafos y los \emph{transformers} como
instancias de una sola construcci\'on geom\'etrica.

\section{C\'omo leer este libro}
\label{sec:how-to-read}

El libro sigue una progresi\'on l\'ogica desde los fundamentos hasta los modelos
avanzados. La \Cref{fig:book-structure} muestra las dependencias entre
cap\'itulos.

\begin{figure}[h]
\centering
\begin{tikzpicture}[
    node distance=0.9cm and 1.4cm,
    box/.style={rectangle, draw, rounded corners, minimum width=2.7cm,
        minimum height=0.8cm, align=center, font=\small},
    found/.style={box, fill=gdlblue!15},
    five/.style={box, fill=gdlgreen!18},
    prac/.style={box, fill=gdlpurple!16},
    arrow/.style={->, thick, >=stealth}
]
\node[found] (intro) {Cap.\ 1\\Introducci\'on};
\node[found, below=of intro] (found) {Cap.\ 2\\Fundamentos};
\node[found, below=of found] (priors) {Cap.\ 3\\Sesgos geom\'etricos};
\node[five, below left=1.1cm and 0.2cm of priors] (grids) {Cap.\ 4--5\\Rejillas, Grupos};
\node[five, below right=1.1cm and 0.2cm of priors] (graphs) {Cap.\ 6--8\\Grafos, Geod\'esicas, Gauges};
\node[prac, below=2.6cm of priors] (apps) {Cap.\ 9--10\\Aplicaciones, Conclusi\'on};
\draw[arrow] (intro) -- (found);
\draw[arrow] (found) -- (priors);
\draw[arrow] (priors) -- (grids);
\draw[arrow] (priors) -- (graphs);
\draw[arrow] (grids) -- (apps);
\draw[arrow] (graphs) -- (apps);
\node[right=1.6cm of priors, font=\footnotesize, align=left] {
    \tikz\fill[gdlblue!15] (0,0) rectangle (0.3,0.3); ~Fundamentos\\[2pt]
    \tikz\fill[gdlgreen!18] (0,0) rectangle (0.3,0.3); ~Las Cinco Ges\\[2pt]
    \tikz\fill[gdlpurple!16] (0,0) rectangle (0.3,0.3); ~Pr\'actica};
\end{tikzpicture}
\caption{Estructura del libro y dependencias entre cap\'itulos.}
\label{fig:book-structure}
\end{figure}

\textbf{Parte I (Fundamentos).} El \Cref{ch:foundations} establece el lenguaje
matem\'atico ---grupos, representaciones, espacios homog\'eneos y variedades---. El
\Cref{ch:priors} introduce la invarianza y la equivarianza y el esquema que
convierte una simetr\'ia en una arquitectura.

\textbf{Parte II (Las Cinco Ges).} Los
\Cref{ch:grids,ch:groups,ch:graphs,ch:geodesics,ch:gauges} desarrollan los cinco
temas por turnos, desde las redes convolucionales sobre rejillas hasta las redes
equivariantes gauge sobre variedades, mostrando c\'omo cada uno extiende al
anterior y c\'omo el paso de mensajes los unifica.

\textbf{Parte III (Pr\'actica).} El \Cref{ch:applications} recorre aplicaciones en
visi\'on por computador, ciencias f\'isicas y datos relacionales, y el
\Cref{ch:conclusion} sintetiza la visi\'on unificada y esboza los problemas
abiertos.

Este es un libro de naturaleza te\'orica, con \'enfasis en el rigor matem\'atico de
sus definiciones, resultados y conexiones formales entre los marcos que presenta.
Se supone que el lector est\'a c\'omodo con el \'algebra lineal, el c\'alculo
multivariable y los fundamentos del aprendizaje autom\'atico; los prerrequisitos
geom\'etricos y de teor\'ia de grupos se desarrollan a medida que se necesitan.
